% Compile with:
%   pdflatex paper_jasss && bibtex paper_jasss && pdflatex paper_jasss && pdflatex paper_jasss
\documentclass[11pt]{article}

\usepackage[a4paper, margin=1in]{geometry}
\usepackage{amsmath,amssymb}
\usepackage{graphicx}
\usepackage{booktabs}
\usepackage{caption}
\usepackage{subcaption}
\usepackage{natbib}
\usepackage{xcolor}
\usepackage[hidelinks]{hyperref}
\usepackage[final,protrusion=true,expansion=false]{microtype}

\graphicspath{%
  {../../out/figures/phase/}%
  {../../out/figures/scenarios/}%
  {../../out/figures/sweeps/}%
  {../../out/figures/policies/}%
}

\title{Three policy targets, three policy levers:\\
\large a stylised agent-based model in which household-sector
inequality, productive-flow share, and population survival are
addressed by distinct fiscal instruments}

\author{[Author]\\
\small (supervision and acknowledgements TBD)}
\date{Draft of \today\ for JASSS}

\begin{document}
\maketitle

\begin{abstract}
We present an agent-based model of an economy with seven occupationally
heterogeneous agent classes and use it to argue that the fiscal-policy
dimensions a wealth-tax debate typically lumps together --- the rate
schedule, the spending package, and the inheritance regime --- act on
distinct outcome metrics and should be evaluated separately.

The model couples production, extraction (rent, interest, capital
return), credit, taxation, inheritance, and a bandwidth-limited
cultural-transmission channel.  We introduce a per-tick scalar
$\Phi_t = F^\text{prod}_t / (F^\text{prod}_t + F^\text{ext}_t)$, the
productive-flow share, defined symmetrically across productive and
extractive accrual.  $\Phi_t$ operationalises the
productive-vs-extractive distinction long named in the heterodox
economic tradition \citep{kauffman1988evolution, mazzucato2018value};
it is not an autocatalytic-set closure in the
\citet{hordijk2004detecting} RAF sense and we are explicit about that.
A Sankey of steady-state flows makes $\Phi$ legible as a diagram.

The principal finding is a factorisation.  A $3 \times 2 \times 3$
factorial over tax schedule, spending package, and inheritance tax
shows that spending raises population alive counts by approximately
40\% without moving $\Phi$ (the spending channels' contribution to
$F^\text{prod}$ is offset, in share terms, by the reduced rentier
extraction it enables); tax-rate schedules compress household-sector
Gini and top-shares roughly fourfold; and inheritance taxation is the
single $\Phi$-most-effective lever, shrinking dynastic rentier stocks
that drive capital-return extraction.  The three levers therefore
address three distinct policy targets and combining them is
multiplicative, not additive.

We extend the factorisation with a fourth lever (an absolute cap
on inheritance) acting on a fourth target (the wealth share held
by the bottom half of the population).  At the joint-calibrated UK
baseline (capital return $r = 0.08$, matching the UK Sunday Times
Rich List Pareto slope), the bottom 50\% holds 0.0\% of household
wealth under status quo and 0.6\% under the Zucman 2\% wealth tax
as proposed.  Under a \pounds 5M absolute cap on inheritance the
bottom-50\% share rises to 19.2\% as a sole intervention and to
25.3\% combined with Zucman 2\%.  Rate-based instruments
redistribute flow; the cap redistributes stock.  The four-lever-
four-target factorisation is the full one.

Thirteen scenarios cluster into three regimes under k-means with
$k=3$ on z-normalised (Gini, $\Phi$, alive); the regime boundaries
are interpretable rather than statistically dictated (silhouette
scores at $k=2$ and $k=3$ are tied at $0.55$).  Eleven real-world
wealth-tax proposals are mapped into the dial panel and compared in
steady state, including the 2024 \citet{zucman2024global} G20
proposal and historical proposals from Piketty, Warren, and Sanders.
The eleven policies separate cleanly on alive, Gini, and top-shares
but not on $\Phi$ at default parameters; whether they separate on
$\Phi$ depends on the assumed rate of capital return.

We provide an ODD-style summary, k-means clustering with reported
silhouette scores in place of a prior ``by-inspection'' assignment,
survival curves underlying the regime classification, and a code and
data release at the project repository.  The cultural-mechanism and
wealth-tax-design results are reported separately in companion
submissions targeting Artificial Life and the Journal of Economic
Interaction and Coordination.
\end{abstract}

%==========================================================================
\section{Introduction}
%==========================================================================

The political case for a wealth tax frequently combines, into a single
proposal, three separable design choices: the rate schedule on wealth
above a threshold, the spending package that the revenue funds, and
the inheritance regime against which both operate.  An agent-based
model of an economy with these three levers as distinct dials can
address whether they are also distinct in effect.  This paper argues
they are.

We approach the question with a stylised stock-flow ABM.  Seven
occupationally typed agents (workers, makers, employers, lenders,
rentiers, speculators, and a single state) exchange monetary state
through production, extraction (rent, interest, capital return),
credit, taxation, inheritance, and a cultural-transmission channel.
Each tick runs the same sequence of phases (Section~\ref{sec:model}).
The model is small enough to read and large enough that the dial
panel exposes the policy choices that the political debate actually
turns on.

We organise the paper around three results.

The first is a measurement device.  We define $\Phi_t$, the
\emph{productive-flow share}, as the share of per-tick monetary
accrual accruing to or being created by productive-class agents.  It
sits in $[0, 1]$ at each tick and characterises whether the
flow regime is dominated by productive transformation or by
extraction.  The Sankey of steady-state per-tick flows is its
diagram (Section~\ref{sec:sankey}).

The second is a regime classification.  Thirteen named scenarios at
$N=100$ cluster into three regions of state space under k-means with
$k=3$: viable autocatalytic, extractive equilibrium, and collapse.
We report the silhouette scores and acknowledge that the
collapse-vs-extractive boundary is the weaker of the two cluster
separations; the choice of three regimes is interpretive, not
dictated by the data.

The third is the factorisation.  A clean factorial over (tax
schedule) $\times$ (spending package) $\times$ (inheritance tax)
shows that each of the three levers moves a different metric.
Spending raises alive counts without moving $\Phi$; tax schedules
compress household-sector Gini and top-shares; inheritance taxation
moves $\Phi$.  The political habit of evaluating a wealth-tax
proposal on a single summary metric collapses three answers into
one.

The model has policy implications that we report only in summary form
here; the wealth-tax-design and calibration findings are the focus of
a companion submission to the Journal of Economic Interaction and
Coordination, and the cultural-mechanism findings are the focus of a
companion submission to Artificial Life.

%==========================================================================
\section{Related work}
\label{sec:related}
%==========================================================================

The closest published economic-ABM neighbour is
\citet{spangler2022distribution}, who couple estate taxation with
binary skill inheritance.  Within the JASSS tradition the model
sits alongside the Sugarscape lineage of
\citet{epstein1996growing} and \citet{wilensky1998wealth}, the
econophysics tradition of \citet{bouchaud2000wealth} and
\citet{yakovenko2009colloquium}, and the kinetic-tax models of
\citet{bertotti2018modanese} and \citet{cardoso2025wealth}.  The
last of these is the closest neighbour on the social-protection
result, namely that what the state does with revenue matters more
than the headline tax rate.  Our factorial-decomposition result
(Section~\ref{sec:factorial}) makes this finding sharper by
separating the spending from the rate schedule and from the
inheritance regime.

For the productive-vs-extractive framing we draw on
\citet{kauffman1988evolution} and \citet{mazzucato2018value} and
distinguish our $\Phi$ from the topological RAF closure of
\citet{hordijk2004detecting} and \citet{hordijk2023emergence}.

%==========================================================================
\section{The model}
\label{sec:model}
%==========================================================================

\subsection{Agent types and population}

Seven agent classes are present at fixed proportions per 100 agents:
72 workers, 9 makers, 7 employers, 5 lenders, 3 rentiers, 3
speculators, and 1 state.  Each agent carries wealth, debt, age,
lifespan drawn from $\mathcal{N}(70, 12)$ truncated to $[20, 120]$,
a five-dimensional skill vector with class-typical biases at
initialisation, peak lifetime wealth, and a bandwidth-at-birth
record.

For scale-up runs we introduce a \texttt{population\_scale}
multiplier that scales agent counts and initial wealth totals
linearly while keeping the state a singleton.  Per-capita quantities
are scale-invariant.

\subsection{Per-tick mechanics}

Each tick runs the following phases in order.

\paragraph{Production.}  Workers are sorted into a queue by skill;
makers and employers hire from the head of the queue.  Each maker
transforms a material input plus paid labour into output at
multiplier $M$, retaining the surplus $\text{output} - \text{inputs}$
as value-add.  Each employer captures a markup on its workers'
produced output; a fraction $0.4$ of gross employer markup is
treated as productive value-add (organising work) and $0.6$ as
zero-sum capture.  The state, when configured, hires a fraction of
unemployed workers from the queue tail and purchases a budgeted
volume of maker output.

\paragraph{Extraction.}  Each worker pays a fixed rent
$\rho \cdot W$ which is distributed across rentiers in proportion to
rentier wealth.  Each rentier earns a capital return $r$ on its
existing wealth.  Each speculator's wealth evolves geometrically with
drift $\mu_s$ and volatility $\sigma_s$.

\paragraph{Consumption and credit.}  Every non-state agent pays
subsistence $\sigma$ per tick.  Agents with insufficient wealth borrow
from the aggregate lender pool, with the borrowed amount allocated to
lenders proportionally to their available capacity.  If aggregate
capacity is exhausted the agent dies (``starves'').

\paragraph{Interest and repayment.}  Indebted agents pay interest at
rate $i$ when liquid; otherwise interest is capitalised.  Agents
beyond a configurable debt ceiling are written off; the loss is
distributed across lenders proportionally to wealth.

\paragraph{Lender deposits and recapitalisation.}  Surplus wealth held
by rentiers, employers, and makers is deposited into the lender pool
at a fixed rate.  The state optionally backstops the lender pool to a
configurable floor.

\paragraph{Tax and UBI.}  Every non-state living agent with positive
wealth pays a stacked-marginal wealth tax
$T(w) = \tau_0 w + \sum_k \tau_k \max(0, w - \theta_k)$.
The state collects $T$ and disburses universal basic income $U$ per
agent to the extent the state's wealth allows.

\paragraph{Life.}  Agents whose age exceeds their lifespan die a
natural death.  A child of the same class is created.  Wealth passes
after an inheritance tax $\iota$ and is capped at
$\bar{w}_\text{inh}$.  The child's skill vector is bandwidth-limited
according to a log-coupling to parent wealth (the cultural channel
documented at length in our ALIFE submission; we summarise it here as
mechanism).

A complete ODD-style summary is in Appendix~\ref{app:odd}.

\subsection{The productive-flow share metric}
\label{sec:phi}

During each tick we accumulate two flow sums treated symmetrically as
productive- and extractive-class accrual.  Productive accrual sums
wages, state transfers, maker value-add (the surplus
$\text{output} - \text{inputs}$), and the productive share of
employer markup.  Extractive accrual sums rent, interest, capital
returns accruing to rentier stocks, and positive speculator returns.
The productive-flow share is
\begin{equation}
\Phi_t \;=\; \frac{F^\text{prod}_t}{F^\text{prod}_t + F^\text{ext}_t}
\;\in\; [0, 1].
\label{eq:phi}
\end{equation}

\paragraph{Nomenclature.}  $\Phi_t$ is \emph{not} an autocatalytic-set
closure in the \citet{hordijk2004detecting} RAF sense; it is a
continuous flow-share that operationalises the
productive-vs-extractive distinction in a measurable form.  The name
``productive-flow share'' is to avoid any implication of a
graph-theoretic property the metric does not claim.

%==========================================================================
\section{The phase portrait and regime classification}
\label{sec:phase}
%==========================================================================

We characterise the model by running thirteen named scenarios at
$N=100$ for 250 ticks across 8 seeds and clustering the endpoints
in $(\text{Gini}, \Phi, \text{alive})$ under k-means with $k=3$
(silhouette $0.55$, with $k=2$ at $0.55$, $k=4$ at $0.42$, $k=5$ at
$0.32$; the data does not statistically prefer $k=3$ over $k=2$ but
the three-regime structure is the more natural reading of the
collapse-vs-extractive distinction).

\begin{figure}[t]
\centering
\begin{subfigure}{0.5\linewidth}
\includegraphics[width=\linewidth]{phase3d_trajectories.png}
\end{subfigure}\hfill
\begin{subfigure}{0.47\linewidth}
\includegraphics[width=\linewidth]{phase3d_shadows.png}
\end{subfigure}
\caption{Phase portrait of thirteen scenarios at $N=100$, 8 seeds.
Three regions are visible: viable autocatalytic ($\Phi \geq 0.4$,
alive $\geq 65$), extractive equilibrium ($\Phi \approx 0.01$,
alive 40--75), collapse ($\Phi = 0$, alive $\leq 20$).
State excluded from Gini and alive count.}
\label{fig:phase}
\end{figure}

\begin{table}[t]
\centering
\small
\caption{Endpoint coordinates by scenario.  State excluded from Gini.}
\label{tab:regimes}
\begin{tabular}{lrrrl}
\toprule
scenario & Gini & $\Phi$ & alive & regime \\
\midrule
socialist             & 0.30 & 0.880 & 99.0 & viable \\
nordic                & 0.30 & 0.701 & 89.9 & viable \\
inheritance\_break\_ubi & 0.57 & 0.667 & 85.9 & viable \\
maker\_economy        & 0.77 & 0.586 & 68.8 & viable \\
tax\_and\_spend       & 0.43 & 0.567 & 83.6 & viable \\
memory\_equalised     & 0.92 & 0.015 & 76.9 & extractive \\
creditless            & 0.87 & 0.013 & 55.2 & extractive \\
current\_world        & 0.89 & 0.011 & 63.1 & extractive \\
memory\_starved       & 0.87 & 0.011 & 55.5 & extractive \\
pure\_speculation     & 0.92 & 0.001 & 57.9 & extractive \\
stevenson\_world      & 0.88 & 0.001 & 20.2 & collapse \\
feudal                & 0.78 & 0.000 & 15.2 & collapse \\
inheritance\_break    & 0.00 & 0.000 &  0.0 & collapse \\
\bottomrule
\end{tabular}
\end{table}

\subsection{Survival curves: showing rather than asserting collapse}

\begin{figure}[t]
\centering
\includegraphics[width=0.97\linewidth]{survival_scenarios.png}
\caption{Survival curves for eight scenarios at $N=100$, 6 seeds.
\textbf{Left:} alive over time.  \textbf{Right:} cumulative excess
mortality.  \texttt{inheritance\_break} loses its population within
30 ticks; \texttt{stevenson\_world} and \texttt{feudal} follow shortly
after.  \texttt{nordic} and \texttt{socialist} hold near full
population.}
\label{fig:survival}
\end{figure}

The collapse regime reaches near-zero alive within 30 ticks for
\texttt{inheritance\_break} (children born with zero wealth starve
immediately) and within 100 ticks for \texttt{stevenson\_world} and
\texttt{feudal} (high rent and capital return extracts faster than
production can replace).  The viable regime is genuinely viable
across the 250-tick run.

\paragraph{The State exclusion.}  Gini and top-share statistics are
reported with the State agent excluded throughout.  Earlier versions
of this work included the State in percentile calculations, with the
result that spend-heavy scenarios reported high Gini because the
State, accumulating tax revenue, dominated the wealth distribution.
The corrected metric shows the spend-heavy scenarios at Gini
$\approx 0.30$, comparable to lived Scandinavian distributions.

%==========================================================================
\section{Visualising the flow regime: the Sankey}
\label{sec:sankey}
%==========================================================================

The Sankey diagram renders mean per-tick wealth flows in steady state
as a network with edge widths proportional to flow magnitude.  Three
policy configurations are shown side by side to make $\Phi$ legible
as a diagram (Figure~\ref{fig:sankey}).

\begin{figure}[t]
\centering
\includegraphics[width=\linewidth]{sankey_flows.png}
\caption{Steady-state wealth flows at $N=10{,}000$, 3 seeds, ticks
80--180.  \textbf{Left (\texttt{uk\_tory\_2010\_24}):} 44\% of
total per-tick flow is the rentier--lender deposit loop; productive
wage flow is 11\%; $\Phi \approx 0.05$.
\textbf{Middle (\texttt{zucman\_2pct\_spend}):} the rentier--lender
loop has shrunk to 14\%; wealth-tax revenue recirculated to workers;
$\Phi \approx 0.44$.
\textbf{Right (\texttt{socialist}):} worker--state--worker
circulation is 40\% of total flow; $\Phi \approx 0.88$.}
\label{fig:sankey}
\end{figure}

The Sankey is the cleanest single artefact for explaining the
productive-flow share to a non-technical audience.  It also reads as
a visualisation of the productive-vs-extractive distinction
independently of the specific $\Phi$ definition.

%==========================================================================
\section{The factorisation: three levers, three targets}
\label{sec:factorial}
%==========================================================================

The principal finding of this paper is that the three policy levers
in our dial panel that the wealth-tax debate typically bundles act on
distinct outcome metrics.  We run a clean factorial over tax
schedule (none, Zucman flat 2\% above threshold $\theta_1$, Zucman
progressive 2/4/8\%), spending package (off, on at UBI 0.4 + state
employment 0.3 + state purchases at scaled rate 5), and inheritance
tax ($\iota \in \{0, 0.4, 0.8\}$).  Eighteen cells, six seeds, 200
ticks.  Table~\ref{tab:factorial} reports marginal means.

\begin{table}[t]
\centering
\small
\caption{Marginal effects of each factor in the
$3 \times 2 \times 3$ factorial, $N=100$, 6 seeds, 200 ticks.}
\label{tab:factorial}
\begin{tabular}{llrrrr}
\toprule
factor & level & alive & Gini & $\Phi$ & top-10 \\
\midrule
tax schedule & none                   & 66.5 & 0.85 & 0.183 & 0.77 \\
             & Zucman flat 2\%        & 67.8 & 0.68 & 0.439 & 0.56 \\
             & Zucman prog 2/4/8\%    & 67.8 & 0.67 & 0.453 & 0.55 \\
\midrule
spending     & off                    & 56.1 & 0.79 & 0.365 & 0.63 \\
             & on                     & 78.6 & 0.71 & 0.351 & 0.62 \\
\midrule
inheritance  & $\iota = 0.0$          & 68.8 & 0.74 & 0.225 & 0.65 \\
             & $\iota = 0.4$          & 67.5 & 0.74 & 0.320 & 0.66 \\
             & $\iota = 0.8$          & 65.8 & 0.74 & 0.529 & 0.56 \\
\bottomrule
\end{tabular}
\end{table}

\paragraph{Spending raises alive without moving $\Phi$.}  Toggling the
spending package on raises the average alive count from 56 to 79,
a 40\% increase; $\Phi$ moves by less than $0.02$ in either
direction (0.365 vs 0.351).  The wage flow added by state spending
is matched, in share terms, by the reduced rentier extraction it
enables.  Spending's headline effect is on mortality, not on the
composition of monetary flow.

\paragraph{Tax schedule compresses household inequality.}  The
Zucman 2\% bracket above threshold moves Gini from 0.85 to 0.68 and
top-10 share from 0.77 to 0.56.  It also raises $\Phi$ from 0.18 to
0.44, because the tax shrinks the rentier stocks that drive the
capital-return component of $F^\text{ext}$.

\paragraph{Inheritance tax is the most $\Phi$-effective lever.}
Moving from $\iota = 0$ to $\iota = 0.8$ raises $\Phi$ from 0.23 to
0.53, a $+0.30$ swing.  Top-10 share moves from 0.65 to 0.56.  The
mechanism is the same as for the wealth tax: estate taxation
truncates dynastic accumulation, shrinking rentier stocks.

\subsection{The (tax\_vs\_ubi) sweep, re-read}

\begin{figure}[t]
\centering
\includegraphics[width=0.97\linewidth]{tax_vs_ubi.png}
\caption{The \texttt{tax\_vs\_ubi} sweep on $(\text{flat\_tax\_rate},
\text{UBI})$.  Eight outcome metrics.  \textbf{Alive} (top left)
peaks at maximum UBI; \textbf{$\Phi$} (top right) is dominated by the
tax-progressivity axis rather than by UBI, consistent with the
factorial above.}
\label{fig:taxubi}
\end{figure}

The factorisation result reads back into the prior 2D sweep
(Figure~\ref{fig:taxubi}): the alive surface is dominated by UBI,
the $\Phi$ surface by the tax axis, and the two objectives are not
co-located.  A policy that maximises alive does not maximise $\Phi$
and vice versa.

%==========================================================================
\section{Wealth-tax policies in summary}
\label{sec:policy}
%==========================================================================

Eleven real-world wealth-tax proposals are encoded as dial settings
in our \texttt{policies.py} library and compared in steady state
(Table~\ref{tab:policies-100}).  Each proposal combines a tax
schedule, a spending package, and an inheritance rate; we run them
as the bundled policy each is typically advertised as.

\begin{table}[t]
\centering
\small
\caption{Eleven real-world wealth-tax proposals at $N=100$, 6 seeds,
250 ticks.  ``Reform UK flavour'' is included as a near-future
political scenario rather than a present policy.}
\label{tab:policies-100}
\begin{tabular}{lrrrrrr}
\toprule
policy & alive & $\Phi$ & top-10 & state $w$ & Gini & metabolic \\
\midrule
uk\_tory\_2010\_24            & 57 & 0.050 & 0.85 & 36{,}541 & 0.90 & 39.0 \\
uk\_labour\_2024              & 57 & 0.230 & 0.75 & 28{,}455 & 0.85 & 39.5 \\
zucman\_2pct\_hoard           & 51 & 0.418 & 0.59 & 17{,}943 & 0.75 & 39.1 \\
zucman\_2pct\_spend           & 80 & 0.412 & 0.54 & 25{,}215 & 0.58 & 43.1 \\
zucman\_progressive\_2\_4\_8 & 80 & 0.420 & 0.54 & 24{,}949 & 0.58 & 43.1 \\
warren\_2020                  & 77 & 0.264 & 0.66 & 34{,}155 & 0.73 & 43.2 \\
sanders\_2020                 & 79 & 0.259 & 0.75 & 33{,}411 & 0.76 & 42.9 \\
piketty\_book                 & 78 & 0.413 & 0.66 & 27{,}550 & 0.72 & 42.9 \\
norway\_actual                & 64 & 0.091 & 0.75 & 40{,}979 & 0.88 & 40.7 \\
spain\_solidarity             & 77 & 0.365 & 0.65 & 28{,}331 & 0.71 & 41.5 \\
reform\_uk\_flavour           & 45 & 0.019 & 0.79 & 28{,}952 & 0.89 & 38.1 \\
\bottomrule
\end{tabular}
\end{table}

Three observations.  First, the eleven policies separate cleanly on
alive (45--80), Gini (0.58--0.90), and top-10 share (0.54--0.85).
Second, the Zucman flat and progressive variants are not separable
at $N=100$ on any metric we measure to two decimal places at default
parameters; whether they separate is a function of the assumed rate
of capital return (the companion JEIC submission shows clean
separation at calibrated parameters).  Third, the
\texttt{reform\_uk\_flavour} prospectus, which is included as a
near-future political projection, gives the lowest alive count (45)
and the lowest $\Phi$ (0.019).

We do not draw \pounds-denominated conclusions from these results in
this paper.  The threshold-comparison and calibration content lives
in the JEIC companion submission.

%==========================================================================
\section{A fourth lever for a fourth target: the bequest cap and
ownership redistribution}
\label{sec:cap-as-fourth}
%==========================================================================

The factorisation in Section~\ref{sec:factorial} reports three
levers (tax schedule, spending package, inheritance taxation)
acting on three targets (Gini, alive, $\Phi$).  None of those three
levers, in any combination we tested, moves a fourth target: the
\emph{ownership share held by the bottom half of the population}.
Under the joint-calibrated UK baseline used in the JEIC companion
submission (capital return $r = 0.08$), the wealth distribution
that matches UK Sunday Times Rich List Pareto slope places the
bottom 50\% at 0.0\% of total household wealth.  Tax schedules,
spending packages, and inheritance taxation each, separately and
in combination, leave this share unchanged or move it by less than
one percentage point.

A fourth lever does change this target.  An absolute cap on
inheritance, set in our companion paper at \pounds 5M based on a
cap-value sweep, raises the bottom-50\% wealth share from 0.0\% to
19.2\% as a sole intervention and to 25.3\% combined with the
Zucman 2\% wealth tax.  Productive-flow share $\Phi$ rises from
near zero under status quo to 0.27--0.37; population survival
improves from 84\% to 91--92\%; the billionaire count under both
$N = 10{,}000$ and the calibrated UK baseline goes to zero or near
zero.  See Table~\ref{tab:cap-fourth}.

\begin{table}[h]
\centering
\small
\caption{Five interventions at the joint-calibrated UK baseline,
extended factorisation.  Bottom-50\% share is the fourth target;
the absolute cap is the fourth lever.}
\label{tab:cap-fourth}
\begin{tabular}{lrrrrr}
\toprule
intervention & alive \% & bottom 50\% share & top 10\% share
& $\Phi$ & $\geq$\pounds 1B \\
\midrule
status quo                & 84\% &  0.0\%  & 100.0\% & 0.00 & 803 \\
Zucman 2\% above \pounds 10M & 91\% &  0.6\%  &  98.9\% & 0.01 & 800 \\
interest ban ($i=0$)      & 89\% &  0.0\%  & 100.0\% & 0.00 & 803 \\
\pounds 5M bequest cap     & 91\% & \textbf{19.2\%} & 54.0\% & 0.27 &   2 \\
cap + Zucman 2\%           & 92\% & \textbf{25.3\%} & 39.9\% & 0.37 &   0 \\
\bottomrule
\end{tabular}
\end{table}

The structural reason a rate-based instrument cannot move the
bottom-50\% share is that rate-based instruments redistribute
\emph{flow} rather than \emph{stock}.  A wealth tax collects
annual revenue from large stocks and the state recirculates it
via spending channels, which improves survival and welfare flows
for the bottom half but does not give them an ownership stake in
the underlying capital.  The cap operates on the stock-transmission
mechanism (inheritance) directly: dynastic wealth above the cap is
either redirected to the state or simply does not propagate to the
next generation, and the wealth that no longer compounds at the
top becomes available, mechanically, to be held more widely.

The methodological observation, given the title of this paper, is
that the three-lever factorisation reported in
Section~\ref{sec:factorial} is incomplete.  There is a fourth
lever (the cap) acting on a fourth target (ownership
redistribution), and the four-lever-four-target factorisation is
the full one.  We report the fourth lever here, refer the
quantitative details to the JEIC companion, and recommend that the
JASSS submission be read as the first three legs of a four-leg
result.

%==========================================================================
\section{Robustness}
\label{sec:robustness}

\subsection{Statistical robustness: 30-seed bootstrap CIs}

We re-ran the eleven-policy comparison at 30 seeds with 95\%
bootstrap confidence intervals on every reported median.  Tighter
confidence intervals confirm the principal qualitative findings
and tighten the equivalence claims.

The flat-vs-progressive Zucman comparison is not separable at 30
seeds: alive median $80$ with $95\%$ CIs $[79, 82]$ for both flat
and progressive; $\Phi$ medians $0.416$ vs $0.422$ with CIs
$[0.402, 0.428]$ vs $[0.413, 0.434]$.  The CIs overlap completely.
The non-separability is rigorous, not seed-noise of an
under-replicated experiment.

The \texttt{reform\_uk\_flavour} alive CI is wide: $[38, 52]$ on a
median of $46$, reflecting the bistable dynamics of the collapse
regime where seeds either die early or grow into a degraded
steady state.

\subsection{One-factor-at-a-time dial sensitivity}

We re-ran \texttt{current\_world} with each principal dial
multiplied by $0.5\times$ and $2\times$ its default in turn.  The
response of each metric to each dial is reported in
Table~\ref{tab:sensitivity}.

\begin{table}[h]
\centering
\small
\caption{OFAT sensitivity at $N=100$, 8 seeds, 250 ticks.  Each
row is a perturbation of one dial; entries are median values.
Baseline median is the all-defaults row.}
\label{tab:sensitivity}
\begin{tabular}{lrrrrr}
\toprule
dial / factor & alive & $\Phi$ & Gini & top-10 \\
\midrule
\emph{baseline}      & 64 & 0.012 & 0.89 & 0.78 \\
wage              $\times 0.5$        & 14 & 0.007 & 0.78 & 0.48 \\
wage              $\times 2.0$        & 59 & 0.014 & 0.91 & 0.96 \\
maker\_multiplier $\times 0.5$        & 50 & 0.008 & 0.94 & 0.99 \\
maker\_multiplier $\times 2.0$        & 66 & 0.023 & 0.88 & 0.78 \\
capital\_return   $\times 0.5$        & 60 & \textbf{0.428} & 0.83 & 0.77 \\
capital\_return   $\times 2.0$        & 66 & 0.000 & 0.93 & 0.89 \\
rent\_share       $\times 0.5$        & 76 & 0.022 & 0.89 & 0.89 \\
rent\_share       $\times 2.0$        & 12 & 0.004 & 0.80 & 0.51 \\
subsistence       $\times 0.5$        & 81 & 0.013 & 0.90 & 0.95 \\
subsistence       $\times 2.0$        & 12 & 0.009 & 0.78 & 0.52 \\
lifespan\_mean    $\times 0.5$        & 72 & 0.013 & 0.90 & 0.88 \\
lifespan\_mean    $\times 2.0$        & 60 & 0.013 & 0.94 & 0.98 \\
bandwidth\_base   $\times 0.5$        & 56 & 0.011 & 0.87 & 0.68 \\
bandwidth\_base   $\times 2.0$        & 78 & 0.015 & 0.92 & 0.98 \\
mutation\_rate    $\times 0.5$        & 66 & 0.013 & 0.89 & 0.79 \\
mutation\_rate    $\times 2.0$        & 62 & 0.012 & 0.89 & 0.78 \\
\bottomrule
\end{tabular}
\end{table}

The qualitative findings.  Survival is dominated by the three
flow-balance dials (wage, rent, subsistence).  $\Phi$ is dominated
by $r$ (halving $r$ raises $\Phi$ from $0.012$ to $0.428$ because
the extractive flow shrinks while productive flow stays roughly
constant).  Inequality metrics respond more uniformly across
dials.  Skill mutation rate is essentially inert at this short
horizon.  The bandwidth-base dial moves alive (by $\pm 14$) and
top-10 share (by $\pm 0.20$) in opposite directions, consistent
with the cultural-economic asymmetry argument elaborated in the
ALIFE companion submission.

\section{Discussion}
%==========================================================================

\subsection{What we believe is novel for a JASSS audience}

The factorisation of policy effects across three distinct levers
(Section~\ref{sec:factorial}, Table~\ref{tab:factorial}) is the
result of this paper.  Each of the three findings (spending raises
alive, tax compresses household inequality, inheritance moves $\Phi$)
is consistent with established economics, but the
\emph{factorisation} --- that they are separable and act on
distinct outcome metrics --- is, to our reading of the kinetic-tax
ABM literature, not previously made explicit.

The productive-flow share $\Phi$ (Section~\ref{sec:phi}) is a
continuous, scalar, per-tick measurement device.  We explicitly
distinguish it from the topological autocatalytic-set closure of
\citet{hordijk2004detecting}.

The Sankey visualisation (Section~\ref{sec:sankey}) is the cleanest
diagram of $\Phi$ as a regime characterisation and we report it as
a methodological contribution in its own right.

\subsection{What the model cannot do}

No spatial structure, no firm formation, fixed agent-type
proportions, toy speculator dynamics that do not produce bubbles, a
single lender pool with no shadow-banking.  The cultural channel,
which we summarise here as mechanism, is reported in full in the
ALIFE companion submission; the wealth-tax design and calibration
content is reported in the JEIC companion submission.

The State exclusion fix and the symmetric $\Phi$ redefinition
changed magnitudes of several headline results between earlier
versions of this work and this one.  None of the qualitative
regime claims changed.

%==========================================================================
\section{Conclusion}
%==========================================================================

A stylised agent-based model of an economy with seven occupationally
typed agents and a dial panel exposing rate, spending, and
inheritance shows that the political habit of evaluating a wealth-tax
proposal on a single summary metric collapses three answers into
one.  Spending keeps people alive; tax rate schedules compress
household-sector inequality; inheritance taxation raises the
productive-flow share.  The three findings are separable and
multiplicative in combination.

A productive-flow share $\Phi$, defined symmetrically across
productive and extractive accrual, makes the third lever
measurable and the regime classification continuous rather than
binary.  The Sankey visualisation makes $\Phi$ legible as a
diagram.

\bibliographystyle{plainnat}
\bibliography{references}

\appendix
\section{ODD-style summary}
\label{app:odd}

\paragraph{Overview.} Seven agent classes (WORKER, MAKER, EMPLOYER,
LENDER, RENTIER, SPECULATOR, STATE) at fixed proportions per 100
agents.  Each agent has wealth, debt, age, lifespan, a five-dim
skill vector, alive flag, peak lifetime wealth, bandwidth-at-birth.

\paragraph{Process.} Per tick: production, extraction, consumption
and credit, interest and repayment, lender deposits, tax and UBI,
life (age, death, reproduction with bandwidth-limited inheritance),
bookkeeping.  Update order is fixed; no inter-tick parallelism.

\paragraph{Stochasticity.}  Per-agent lifespan and initial wealth
share, debt at birth for workers, speculator returns, skill
mutation, and seed-controlled initialisation are stochastic.  All
other quantities are deterministic given dials.

\paragraph{Observation.}  Metrics reported per tick: Gini
(state-excluded), $\Phi$, alive, total wealth, total debt, cultural
breadth, cultural diversity, knowledge lost, births and deaths
counters, top-1 and top-10 shares (raw and mortality-corrected),
state wealth, wealth by type.

\paragraph{Code.}  \texttt{sim.py} for the core model;
\texttt{scenarios.py} for the 13-scenario library;
\texttt{policies.py} for the 11 real-world policy library;
\texttt{factorial.py} for the factorial reported here.  Default
seed 42, default \texttt{population\_scale} 1.

\section{Dial reference}
\label{app:dials}

The dial panel and default values follow.  Tax thresholds are in
model units; for the JEIC companion submission we adopt the
convention 1 unit = \pounds 100k for political-policy scenarios.

\begin{table}[h]
\centering
\small
\begin{tabular}{lrl}
\toprule
dial & default & group \\
\midrule
wage                          & 1.0 & production \\
maker\_multiplier             & 1.6 & production \\
employer\_markup              & 0.4 & production \\
capital\_return ($r$)         & 0.04 & capital \\
rent\_share\_of\_wage         & 0.35 & capital \\
interest\_rate ($i$)          & 0.06 & capital \\
flat\_tax\_rate               & 0.0  & state \\
ubi                           & 0.0  & state \\
state\_employs\_fraction      & 0.0  & state \\
state\_buys\_maker\_output    & 0.0  & state \\
subsistence                   & 0.6  & life \\
debt\_ceiling                 & 25   & life \\
lifespan\_mean / std          & 70 / 12 & life \\
inheritance\_cap              & $\infty$ & transmission \\
inheritance\_tax\_rate        & 0.0   & transmission \\
memory\_bandwidth\_base ($b_0$) & 0.5 & transmission \\
memory\_bandwidth\_wealth\_scale ($b_w$) & 0.5 & transmission \\
skill\_mutation\_rate         & 0.05 & transmission \\
population\_scale             & 1    & logistics \\
seed                          & 42   & logistics \\
\bottomrule
\end{tabular}
\end{table}

\end{document}
